\section*{第 \ref{ch:g9:metals-acids-corrosion} 章——金属：与酸的反应和腐蚀}

\begin{solution}{exo:g9:metals-acids-corrosion:1}
氢气：将点燃的木条靠近试管口，发出爆鸣声。
\end{solution}

\begin{solution}{exo:g9:metals-acids-corrosion:2}
铁、锌和铝能发生反应；铜和金则不能。
\end{solution}

\begin{solution}{exo:g9:metals-acids-corrosion:3}
氧气和水两者都需要。
\end{solution}

\begin{solution}{exo:g9:metals-acids-corrosion:4}
反应中存在但不参与反应的\omterm{def:g9:ions:ion}{离子}。以
盐酸为例：氯\omterm{def:g9:ions:ion}{离子} \ce{Cl-}。
\end{solution}

\begin{solution}{exo:g9:metals-acids-corrosion:5}
\ce{Mg(s) + 2H+(aq) -> Mg^{2+}(aq) + H2(g)}；反应前后两侧电荷均为 $+2$。
\end{solution}

\begin{solution}{exo:g9:metals-acids-corrosion:6}
加入氢氧化钠，生成绿色\omterm{def:g9:ions:precipitate}{沉淀} \ce{Fe(OH)2}，
说明存在亚铁\omterm{def:g9:ions:ion}{离子}。
\end{solution}

\begin{solution}{exo:g9:metals-acids-corrosion:7}
煮沸可除去水中\omterm{def:g2:mixing-and-dissolving:dissolve}{溶解}的\omterm{def:g4:air-a-mixture-of-gases:air}{空气}，油层则阻止\omterm{def:g4:air-a-mixture-of-gases:air}{空气}
再次溶入水中：铁钉处有水，却没有氧气。
\end{solution}

\begin{solution}{exo:g9:metals-acids-corrosion:8}
\omterm{def:g2:mixing-and-dissolving:dissolve}{溶解}在水中的盐溅到汽车上，会加快生锈。
\end{solution}

\begin{solution}{exo:g9:metals-acids-corrosion:9}
形成的氧化层附着在金属上，将金属与\omterm{def:g4:air-a-mixture-of-gases:air}{空气}隔开：
\omterm{def:g9:metals-acids-corrosion:corrosion}{腐蚀}只停留在表面。
\end{solution}

\begin{solution}{exo:g9:metals-acids-corrosion:10}
涂漆、涂油或涂脂，或\omterm{def:g9:metals-acids-corrosion:galvanising}{镀锌}（或用不锈钢制造）。
\end{solution}

\begin{solution}{exo:g9:metals-acids-corrosion:11}
\ce{2Al(s) + 6H+(aq) -> 2Al^{3+}(aq) + 3H2(g)}：两侧均有 2 个 Al 和 6 个 H，
电荷均为 $+6$。对 200 个 Al \omterm{def:g7:atoms-and-molecules:atom}{原子}而言：$200 \times \frac{3}{2} = 300$
个氢气\omterm{def:g7:atoms-and-molecules:molecule}{分子}。
\end{solution}

\begin{solution}{exo:g9:metals-acids-corrosion:12}
划痕周围的锌先于钢被\omterm{def:g9:metals-acids-corrosion:corrosion}{腐蚀}，
只要附近仍有锌，钢便得以保护。
\end{solution}

\begin{solution}{pb:g9:metals-acids-corrosion:1}
\textbf{1.} \ce{Zn(s) + 2H+(aq) -> Zn^{2+}(aq) + H2(g)}。

\textbf{2.} 氢气：点燃的木条靠近时发出爆鸣声。

\textbf{3.} 加入氢氧化钠后生成白色\omterm{def:g9:ions:precipitate}{沉淀}
\ce{Zn(OH)2}。

\textbf{4.} 氯\omterm{def:g9:ions:ion}{离子} \ce{Cl-}。

\textbf{5.} $125.6 - 113.5 = \qty{12.1}{g}$。

\textbf{6.} 只要锌与酸反应，就持续放出氢气；
锌耗尽后，气泡便停止产生。

\textbf{7.} 钢中的铁也会反应，生成氢气和
\ce{Fe^{2+}} \omterm{def:g9:ions:ion}{离子}；加入氢氧化钠后会生成绿色\omterm{def:g9:ions:precipitate}{沉淀}。

\textbf{8.} 锌先于钢被\omterm{def:g9:metals-acids-corrosion:corrosion}{腐蚀}，使钢
与水和\omterm{def:g4:air-a-mixture-of-gases:air}{空气}隔绝。

\textbf{9.} $2 \times 10.0 \times 10.0 = \qty{200}{cm^2}$。

\textbf{10.} $12.1 \div 7.13 \approx \qty{1.70}{cm^3}$。

\textbf{11.} $1.70 \div 200 = \qty{0.0085}{cm}$。

\textbf{12.} $0.0085 \times 10\,000 = \qty{85}{\micro m}$：锌层
厚约 \qty{85}{\micro m}。
\end{solution}
